\newpage
\section{Motivação}
\label{sc:motivacao}


A justificativa para a realização deste estudo fundamenta-se nos severos desafios de engenharia e infraestrutura impostos pela operação de Modelos de Linguagem de Grande Porte (LLMs) em cenários de produção. Embora o surgimento da arquitetura de Geração Aumentada de Recuperação (RAG) tenha mitigado o problema das alucinações ao fornecer dados de fontes externas em tempo real, sua viabilidade comercial enfrenta uma barreira crítica: a ineficiência no consumo de recursos. À medida que o volume de dados organizacionais cresce, o processo tradicional de recuperar e injetar fragmentos extensos de texto diretamente nos prompts resulta em um aumento exponencial no consumo de tokens, saturação da janela de contexto e elevação dos custos com computação em nuvem. No ecossistema tecnológico atual, há uma lacuna evidente na compreensão de como as diferentes abordagens de manipulação desse contexto — como os métodos Stuff, Refine, Map Reduce, Map Re-rank, Query Step-Down e Reciprocal RAG \cite{topsakal2023langchain} — afetam diretamente o uso de hardware, especificamente o consumo de CPU, memória RAM e tempo de resposta (latência). Investigar esses parâmetros é fundamental para a Ciência da Computação e Engenharia de Software, pois o sucesso de aplicações de inteligência artificial depende intrinsecamente da habilidade de equilibrar a acurácia da resposta gerada com a sustentabilidade financeira e computacional da infraestrutura.

A relevância deste tema acentua-se diante das tendências de mercado e dos recentes avanços científicos na área de Processamento de Linguagem Natural (PLN). Estudos recentes de otimização sistemática por varredura em grade (grid-search), que chegam a realizar mais de 23.000 iterações de testes cruzados em domínios complexos (como relatórios financeiros e dados médicos) \cite{sakar2025rag}, apontam que a escolha cega de componentes pode inviabilizar projetos de IA. A compressão de contexto e o uso estratégico de filtros de similaridade surgem não apenas como alternativas de refinamento, mas como requisitos obrigatórios para reduzir o tráfego desnecessário de dados. Consequentemente, a motivação para este trabalho decorre da necessidade prática de mapear esses cenários e propor um guia metodológico claro. Ao correlacionar a eficiência algorítmica com o impacto físico no hardware, este estudo conecta-se a um objetivo muito maior: democratizar e viabilizar a implementação de sistemas RAG robustos e escaláveis, reduzindo o desperdício computacional e otimizando a entrega de valor tecnológico tanto no ambiente acadêmico quanto na indústria de software.